\documentclass[a4paper, 12pt]{article}

\usepackage[utf8]{inputenc}
\usepackage[T1]{fontenc}
\usepackage{textcomp}
\usepackage{amssymb}
\usepackage{newtxtext} \usepackage{newtxmath}
\usepackage{amsmath, amssymb}
\newtheorem{problem}{Problem}
\newtheorem{example}{Example}
\newtheorem{lemma}{Lemma}
\newtheorem{theorem}{Theorem}
\newtheorem{problem}{Problem}
\newtheorem{example}{Example} \newtheorem{definition}{Definition}
\newtheorem{lemma}{Lemma}
\newtheorem{theorem}{Theorem}
\DeclareMathAlphabet{\mathcal}{OMS}{cmsy}{m}{n}

\begin{document}

La bifurcación que se abre ante mí es la siguiente: o bien la vida no tiene
sentido, en cuyo caso no veo otro camino que la muerte voluntaria; o bien la
vida lo tiene, en cuyo caso dicho sentido determinará el camino. Me rehuso a
conceder a la vida un sentido sobrenatural: la vida es parte de la naturaleza.
El único otro tipo de sentido que puede dársele es ético. Por lo que mi dilema
se reduce al siguiente: o logro dotar a mi vida de un sentido ético, o abrazo la
muerte.

~ 

Uno podría decir: ¿por qué morir, si estás avanzando en la vida? Pero todo lo
que he logrado es vano. Hace cuatro años lo hubiera dado todo por, digamos,
publicar un trabajo propio en un journal de \textit{Nature}, por trabajar en
ciencia, por recorrer el mundo, pero habiéndolo alcanzado me digo a mí mismo, ¿y
qué? Frente a mí hay una silla, una mesa, una almohada: ¿qué me distingue de
estas cosas? En mi vida romántica o bien disimulo esta forma de sentirme o, al
mostrarla, soy del todo incomprendido. Las personas responden con extrema
facilidad a la pregunta de cuál es el sentido de la vida, como si preguntáramos
cuál es el color del cielo. La respuesta de la filosofía analítica, que nos dice
que la pregunta no tiene sentido, no ayuda en nada. Las sentencias de la ética,
como Wittgenstein mostró irrefutablemente, no se distinguen en nada de otras
proposiciones. Según él existe algo que nos permite trascender lo proposicional,
algo que nos permite ir más allá del mundo. ¿Pero qué es?

~ 

Esta mañana desperté temprano y, antes de trabajar, fui a dar de comer a los
pájaros en un parque y a leer las enseñanzas de Cristo. En verdad, fue el más
sabio de los hombres. Me parece a mí que \textit{si es que} existe un sentido
ético de la vida, debe ser el de la ética cristiana. A su vez, cuando me siento
solo, que es la mayor parte del tiempo, siento algo así como «Dios te está
llamando» ---aunque esas palabras son inexactas---. Me parece a mí que es
verdadero que «somos la luz del mundo»: en la tiniebla universal arde la luz de
la consciencia. Pero estoy lleno de sombras y la línea que divide al mal del
bien delimita mi vida.

~ 

Desde un punto de vista psicológico, soy consciente que existe una tristeza
subyaciente a esta incertidumbre, a este deseo de descubrir si la vida tiene
sentido y de morir si no lo tiene. Soy consciente, además, de que esta tristeza,
que es como el manantial de todas estas inquietudes, está directamente atada al
suicidio de Paulina. Por predecible que fuera, no estaba preparado para perder
de esa manera a una mujer que fue por tantos años mi amante, mi amiga, mi
compañera: lejos o cerca, era todas estas cosas. Desde que consumé mi duelo,
algo que sucedió más de un año después de su muerte, toda en mí es un agua
removida. Hundo las manos y trato de asir algo, como el ciego que busca con las
manos una perla a plena luz del día.























\end{document}



